\section{Implementation Details}
\label{sec:implementation}

\subsection{Development Environment}

Our implementation targets:
\begin{itemize}
    \item \textbf{NS-3 version}: ns-3.40+ (tested on ns-3-dev)
    \item \textbf{MPI library}: OpenMPI 4.1+ or MPICH 3.4+
    \item \textbf{Compiler}: GCC 9+ or Clang 10+ with C++17 support
    \item \textbf{Build system}: CMake 3.18+ (NS-3's default)
    \item \textbf{OS}: Linux (Ubuntu 20.04/22.04, CentOS 8+)
\end{itemize}

\subsection{Build Configuration}

MPI support must be enabled during NS-3 configuration:

\begin{lstlisting}[style=bash, caption={NS-3 MPI build configuration}]
# Configure NS-3 with MPI support
./ns3 configure --enable-mpi --enable-examples --enable-tests

# Verify MPI is enabled
./ns3 show config | grep MPI
# Should output: MPI: enabled

# Build with parallelism
./ns3 build -j$(nproc)
\end{lstlisting}

The build system automatically detects MPI using CMake's \texttt{FindMPI} module. If multiple MPI implementations exist, specify explicitly:

\begin{lstlisting}[style=bash]
# Use specific MPI implementation
export MPI_HOME=/usr/lib/x86_64-linux-gnu/openmpi
./ns3 configure --enable-mpi
\end{lstlisting}

\subsection{Packet Serialization}

\subsubsection{WiFi Metadata Preservation}

Transmitting WiFi packets via MPI requires serializing not just packet data but also WiFi-specific metadata. NS-3's \texttt{WifiPpdu} contains:

\begin{itemize}
    \item Physical layer parameters (modulation, coding rate)
    \item Transmission duration
    \item MPDU boundaries (for A-MPDU aggregation)
    \item Scrambling initialization
\end{itemize}

Our serialization preserves essential information:

\begin{lstlisting}[style=cpp, caption={WiFi packet serialization structure}]
struct WifiTxMetadata {
    // Source identification
    uint32_t srcRank;
    uint32_t srcDeviceId;
    
    // Physical parameters
    double txPowerDbm;
    uint32_t channelNumber;
    uint16_t channelWidth;  // MHz
    
    // Position (for propagation calculation)
    double posX, posY, posZ;
    
    // Timing
    uint64_t txStartTime;  // Nanoseconds
    uint32_t txDurationUs; // Microseconds
    
    // Packet data
    uint32_t packetSize;
    uint8_t modulation;    // WifiModulationClass enum
    uint8_t preamble;      // WifiPreamble enum
    
    // Followed by actual packet bytes
};
\end{lstlisting}

\subsubsection{Efficient Binary Encoding}

For large simulations with thousands of packets/second, serialization efficiency matters. We use binary encoding:

\begin{lstlisting}[style=cpp, caption={Binary packet serialization}]
std::vector<uint8_t> SerializeWifiPacket(
    Ptr<const WifiPpdu> ppdu,
    const WifiTxMetadata& metadata)
{
    // Calculate total size
    size_t totalSize = sizeof(WifiTxMetadata) + 
                       metadata.packetSize;
    
    std::vector<uint8_t> buffer(totalSize);
    
    // Copy metadata
    std::memcpy(buffer.data(), &metadata, 
                sizeof(WifiTxMetadata));
    
    // Extract and copy packet data
    Ptr<Packet> packet = ppdu->GetPacket()->Copy();
    packet->CopyData(buffer.data() + 
                     sizeof(WifiTxMetadata),
                     metadata.packetSize);
    
    return buffer;
}
\end{lstlisting}

This avoids text-based formats (XML, JSON) which would add significant overhead.

\subsection{Position Management and Mobility}

\subsubsection{Static Position Handling}

For stationary devices, positions are sent once during device registration:

\begin{lstlisting}[style=cpp, caption={Static position registration}]
void RemoteYansWifiChannelStub::RegisterDevice(
    Ptr<YansWifiPhy> phy)
{
    Ptr<MobilityModel> mobility = phy->GetMobility();
    Vector pos = mobility->GetPosition();
    
    DeviceRegMessage msg;
    msg.deviceId = m_deviceCounter++;
    msg.posX = pos.x;
    msg.posY = pos.y;
    msg.posZ = pos.z;
    msg.isStatic = true; // No future updates needed
    
    MPI_Send(&msg, sizeof(msg), MPI_BYTE, 
             m_channelRank, MSG_DEVICE_REG, 
             MPI_COMM_WORLD);
}
\end{lstlisting}

\subsubsection{Mobile Device Position Updates}

For mobile devices, positions must be updated as they move. Two strategies exist:

\paragraph{Periodic Updates}
Send position updates at fixed intervals:

\begin{lstlisting}[style=cpp, caption={Periodic position updates}]
void RemoteYansWifiChannelStub::SchedulePositionUpdate()
{
    static const Time updateInterval = MilliSeconds(100);
    
    for (auto& [id, phy] : m_phyMap) {
        Ptr<MobilityModel> mobility = phy->GetMobility();
        Vector newPos = mobility->GetPosition();
        
        // Send update to rank 0
        PositionUpdateMessage msg;
        msg.deviceId = id;
        msg.posX = newPos.x;
        msg.posY = newPos.y;
        msg.posZ = newPos.z;
        msg.timestamp = Simulator::Now().GetNanoSeconds();
        
        MPI_Send(&msg, sizeof(msg), MPI_BYTE,
                 m_channelRank, MSG_POSITION_UPDATE,
                 MPI_COMM_WORLD);
    }
    
    Simulator::Schedule(updateInterval, 
        &RemoteYansWifiChannelStub::SchedulePositionUpdate,
        this);
}
\end{lstlisting}

\paragraph{Event-Driven Updates (Future Work)}
Hook into NS-3's \texttt{CourseChange} trace to send updates only when position changes significantly:

\begin{lstlisting}[style=cpp, caption={Course change callback (planned)}]
void RemoteYansWifiChannelStub::OnCourseChange(
    Ptr<const MobilityModel> mobility)
{
    Vector newPos = mobility->GetPosition();
    Vector oldPos = m_lastPositions[mobility];
    
    // Send update only if moved > threshold
    double distance = CalculateDistance(newPos, oldPos);
    if (distance > m_updateThreshold) {
        SendPositionUpdate(mobility, newPos);
        m_lastPositions[mobility] = newPos;
    }
}
\end{lstlisting}

\textbf{Current Status}: Our implementation uses static positions. Dynamic mobility support is documented as high-priority future work (Section~\ref{sec:future}).

\subsection{Propagation Model Integration}

\subsubsection{Two-Ray Ground Model for V2X}

Vehicular scenarios use the two-ray ground propagation model:

\begin{equation}
P_r = \frac{P_t G_t G_r h_t^2 h_r^2}{d^4 L}
\end{equation}

where:
\begin{itemize}
    \item $P_r$ = received power
    \item $P_t$ = transmitted power
    \item $G_t, G_r$ = antenna gains (typically 1)
    \item $h_t, h_r$ = antenna heights (vehicles: 1.5m, RSU: 5m)
    \item $d$ = distance between transmitter and receiver
    \item $L$ = system loss factor
\end{itemize}

\begin{lstlisting}[style=cpp, caption={Two-ray ground model configuration}]
Ptr<TwoRayGroundPropagationLossModel> twoRay = 
    CreateObject<TwoRayGroundPropagationLossModel>();
twoRay->SetAttribute("HeightAboveZ", DoubleValue(1.5));
// Frequency and wavelength set automatically by WiFi PHY
processor->SetPropagationLossModel(twoRay);
\end{lstlisting}

\subsubsection{Log-Distance Model for Indoor}

Home WiFi scenarios use log-distance path loss:

\begin{equation}
PL(d) = PL(d_0) + 10 n \log_{10}\left(\frac{d}{d_0}\right) + X_\sigma
\end{equation}

where:
\begin{itemize}
    \item $PL(d_0)$ = reference path loss at $d_0 = 1$m (40-50 dB for indoor)
    \item $n$ = path loss exponent (2.7-3.5 for indoor environments)
    \item $X_\sigma$ = Gaussian random variable (shadowing)
\end{itemize}

\begin{lstlisting}[style=cpp, caption={Log-distance model for indoor WiFi}]
Ptr<LogDistancePropagationLossModel> logDistance = 
    CreateObject<LogDistancePropagationLossModel>();
logDistance->SetAttribute("Exponent", DoubleValue(3.0));
logDistance->SetAttribute("ReferenceDistance", 
                          DoubleValue(1.0));
logDistance->SetAttribute("ReferenceLoss", 
                          DoubleValue(40.0));
processor->SetPropagationLossModel(logDistance);
\end{lstlisting}

\subsubsection{Propagation Delay}

Constant-speed propagation delay:

\begin{equation}
\tau = \frac{d}{c}
\end{equation}

where $c = 3 \times 10^8$ m/s (speed of light). For typical WiFi ranges (< 1km), delays are sub-microsecond but must be modeled for accurate MAC timing.

\subsection{Code Organization}

\begin{table}[htbp]
\centering
\caption{Implementation file structure}
\label{tab:file-structure}
\small
\begin{tabular}{@{}lp{8cm}@{}}
\toprule
\textbf{File/Directory} & \textbf{Purpose} \\
\midrule
\texttt{src/wifi/model/} & \\
\quad\texttt{remote-yans-wifi-channel-stub.h/cc} & MPI stub implementation \\
\quad\texttt{wifi-channel-mpi-processor.h/cc} & Channel processor for rank 0 \\
\quad\texttt{wifi-mpi-messages.h} & MPI message structures \\
\midrule
\texttt{src/wifi/helper/} & \\
\quad\texttt{mpi-wifi-helper.h/cc} & Configuration helper (planned) \\
\midrule
\texttt{scratch/} & \\
\quad\texttt{test.cc} & Basic MPI integration test \\
\quad\texttt{home-wifi-scenario.cc} & 8-device home WiFi \\
\quad\texttt{home-wifi-with-pcap.cc} & Hybrid MPI + PCAP \\
\quad\texttt{v2x-mpi-scenario.cc} & Vehicular communication \\
\midrule
\texttt{doc/mpi-wifi-report/} & \\
\quad\texttt{main.tex} & This report \\
\quad\texttt{sections/*.tex} & Report sections \\
\quad\texttt{figures/} & Diagrams and plots \\
\bottomrule
\end{tabular}
\end{table}

\subsection{Error Handling and Logging}

\subsubsection{MPI Error Handling}

MPI operations can fail; we check return codes:

\begin{lstlisting}[style=cpp, caption={MPI error handling}]
int SendTxRequest(const TxRequestMsg& msg)
{
    int ret = MPI_Send(&msg, sizeof(msg), MPI_BYTE,
                       m_channelRank, MSG_TX_REQUEST,
                       MPI_COMM_WORLD);
    
    if (ret != MPI_SUCCESS) {
        NS_FATAL_ERROR("MPI_Send failed: rank " 
                      << m_localRank 
                      << ", error code " << ret);
    }
    
    return ret;
}
\end{lstlisting}

\subsubsection{Comprehensive Logging}

NS-3's logging system provides multi-level output:

\begin{lstlisting}[style=cpp, caption={Logging levels for debugging}]
// Enable component logging
LogComponentEnable("RemoteYansWifiChannelStub", 
                   LOG_LEVEL_INFO);
LogComponentEnable("WifiChannelMpiProcessor", 
                   LOG_LEVEL_DEBUG);

// Add context to log messages
LogComponentEnableAll(LOG_PREFIX_TIME);   // Timestamp
LogComponentEnableAll(LOG_PREFIX_NODE);   // Node ID
LogComponentEnableAll(LOG_PREFIX_FUNC);   // Function name
\end{lstlisting}

Example output:
\begin{verbatim}
+2.008050000s -1 RemoteYansWifiChannelStub:Send() 
    Device 0 transmitted via MPI, Power: 20.0 dBm
+2.008051000s  0 WifiChannelMpiProcessor:ProcessTx()
    Processing TX from rank 1, device 0
+2.008052000s  0 WifiChannelMpiProcessor:SendRxNotify()
    -> Receiver rank 2, device 1: RX power -45.2 dBm
\end{verbatim}

\subsection{Configuration and Usage}

\subsubsection{Example Scenario Setup}

A complete MPI WiFi scenario requires:

\begin{lstlisting}[style=cpp, caption={Minimal MPI WiFi setup}]
#include "ns3/mpi-interface.h"
#include "ns3/remote-yans-wifi-channel-stub.h"
#include "ns3/wifi-channel-mpi-processor.h"

int main(int argc, char* argv[])
{
    // Initialize MPI
    MpiInterface::Enable(&argc, &argv);
    uint32_t rank = MpiInterface::GetSystemId();
    uint32_t size = MpiInterface::GetSize();
    
    if (rank == 0) {
        // Rank 0: Channel processor
        auto processor = 
            CreateObject<WifiChannelMpiProcessor>();
        // Configure propagation models...
    } else {
        // Ranks 1-N: Devices
        NodeContainer nodes;
        nodes.Create(devicesPerRank);
        
        // Create MPI channel stub
        auto channelStub = 
            CreateObject<RemoteYansWifiChannelStub>();
        channelStub->SetRemoteChannelRank(0);
        channelStub->SetLocalDeviceRank(rank);
        
        // Install WiFi...
    }
    
    Simulator::Run();
    Simulator::Destroy();
    MpiInterface::Disable();
}
\end{lstlisting}

\subsubsection{Command-Line Parameters}

Our scenarios support extensive configuration:

\begin{lstlisting}[style=bash, caption={Command-line configuration}]
# Home WiFi with custom parameters
mpirun -np 4 ./ns3 run "scratch/home-wifi-scenario \
    --simTime=600 \
    --payloadSize=1400 \
    --enablePcap=true"

# V2X with more vehicles
mpirun -np 6 ./ns3 run "scratch/v2x-mpi-scenario \
    --lanes=6 \
    --vehiclesPerLane=100 \
    --simTime=900 \
    --vehicleSpeed=35"
\end{lstlisting}

\subsection{Memory Management}

\subsubsection{Packet Buffer Management}

MPI message buffers require careful management:

\begin{lstlisting}[style=cpp, caption={Buffer lifecycle management}]
void SendPacket(Ptr<const Packet> packet)
{
    uint32_t size = packet->GetSize();
    
    // Allocate buffer
    uint8_t* buffer = new uint8_t[size];
    packet->CopyData(buffer, size);
    
    // Send via MPI
    MPI_Send(buffer, size, MPI_BYTE, 
             destRank, tag, MPI_COMM_WORLD);
    
    // Free immediately (MPI copies data)
    delete[] buffer;
}
\end{lstlisting}

\subsubsection{Device Registry}

The channel processor maintains device state efficiently:

\begin{lstlisting}[style=cpp, caption={Efficient device registry}]
// Use pair<rank,deviceId> as key for O(log n) lookup
std::map<std::pair<uint32_t, uint32_t>, DeviceInfo> 
    m_deviceRegistry;

// Reserve capacity for expected device count
m_deviceRegistry.reserve(expectedDeviceCount);
\end{lstlisting}

\subsection{Performance Optimization Strategies}

\subsubsection{Message Batching (Future)}

Currently, each transmission sends separate MPI messages. Future work could batch multiple transmissions:

\begin{lstlisting}[style=cpp, caption={Message batching concept}]
struct BatchedTxRequest {
    uint32_t count;
    TxRequestMsg requests[MAX_BATCH];
};

// Collect transmissions for a time window
// Send batch via single MPI_Send()
\end{lstlisting}

\subsubsection{Non-Blocking MPI (Future)}

Asynchronous MPI could hide communication latency:

\begin{lstlisting}[style=cpp, caption={Non-blocking MPI concept}]
MPI_Request request;
MPI_Isend(&msg, sizeof(msg), MPI_BYTE, 
          destRank, tag, MPI_COMM_WORLD, &request);
// Continue processing...
MPI_Wait(&request, MPI_STATUS_IGNORE);
\end{lstlisting}

Both optimizations are documented for future implementation.
